\lecture{8.8}{Doubly Robust Longitudinal Estimation}
\label{lec:8-8}

\subsection{Longitudinal AIPTW}

Longitudinal augmented inverse probability weighting combines the sequential
outcome-regression representation with the longitudinal weighting
representation. For the always-treated regime, define
\[
  \gbar_{j,n}(\Lbar_j)=\prod_{m=0}^j g_{m,n}(\Lbar_m),
\]
and let $\Qbar_{j,n}$ estimate the sequential outcome regressions from
\cref{lec:8-3}. The baseline plug-in part is the sequential g-computation
estimate
\[
  \frac{1}{n}\sum_{i=1}^n \Qbar_{1,n}(L_{0i}).
\]
The weighted residual corrections are applied after this baseline averaging. For
$j=0,\ldots,K-1$, define
\[
  \Delta_{ji}
  =
  \Qbar_{j+2,n}(\Lbar_{j+1,i})
  -
  \Qbar_{j+1,n}(\Lbar_{j,i}),
\]
and define the final residual as
\[
  \Delta_{Ki}
  =
  Y_i-\Qbar_{K+1,n}(\Lbar_{K,i}).
\]
The longitudinal AIPTW estimator can then be written as
\[
  \psi_{n,\AIPTW}
  =
  \frac{1}{n}\sum_{i=1}^n \Qbar_{1,n}(L_{0i})
  +
  \frac{1}{n}\sum_{i=1}^n
  \sum_{j=0}^K
  \frac{A_{0i}\cdots A_{ji}}{\gbar_{j,n}(\Lbar_{j,i})}
  \Delta_{ji},
\]
with the convention that $\Qbar_{K+1,n}(\Lbar_K)$ is the final outcome
regression. The baseline plug-in term is not weighted; only the residual
corrections are weighted by regime-following inverse probabilities.

For $K=1$ and regime $(1,1)$, this becomes
\[
\begin{aligned}
  \psi_{n,\AIPTW}
  &=
  \frac{1}{n}\sum_{i=1}^n \Qbar_{1,n}(L_{0i})\\
  &\quad+
  \frac{1}{n}\sum_{i=1}^n
  \frac{A_{0i}}{g_{0,n}(L_{0i})}
  \{\Qbar_{2,n}(L_{0i},L_{1i})-\Qbar_{1,n}(L_{0i})\}\\
  &\quad+
  \frac{1}{n}\sum_{i=1}^n
  \frac{A_{0i}A_{1i}}{g_{0,n}(L_{0i})g_{1,n}(\Lbar_{1i})}
  \{Y_i-\Qbar_{2,n}(L_{0i},L_{1i})\}.
\end{aligned}
\]
The first line is sequential g-computation. The second and third lines are
weighted residual corrections at time 0 and time 1.

\subsection{Nuisance Functions Over Time}

Longitudinal doubly robust estimation uses two nuisance sequences:
\[
  \Qbar=(\Qbar_1,\ldots,\Qbar_{K+1}),
  \qquad
  g=(g_0,\ldots,g_K).
\]
The $\Qbar$ sequence describes outcome and pseudo-outcome regressions. The $g$
sequence describes the treatment process. The estimator uses both because each
sequence can help correct errors in the other.

\subsection{Double Robustness in Multi-Timepoint Settings}

\begin{definition}[Longitudinal double robustness]
A longitudinal estimator is doubly robust if it consistently estimates
$\E\{Y(\abar)\}$ when either the required outcome-regression sequence is
consistently estimated or the required treatment-mechanism sequence is
consistently estimated, under the other identification assumptions.
\end{definition}

For multiple time points, there are also mixed robustness patterns. Roughly, the
estimator can remain consistent when the treatment mechanism is consistently
estimated for earlier parts of the process and the outcome regressions are
consistently estimated for later parts. The exact statement depends on the
recursive representation of the estimator, but the practical message is that
longitudinal AIPTW has more robustness than pure g-computation or pure IPTW.
Mixed robustness is not an arbitrary mix-and-match property: the exact
combinations of correctly estimated $\Qbar$- and $g$-components that yield
consistency depend on the recursive estimator and its indexing.

\subsection{Correct $\Qbar$-Sequence}

If the sequential outcome regressions are consistently estimated, then the
g-computation component already converges to the target:
\[
  \frac{1}{n}\sum_{i=1}^n\Qbar_{1,n}(L_{0i})
  \rightarrow \E\{Y(\abar)\}.
\]
The weighted residual augmentation terms have mean approximately zero because
the residuals
\[
  \Delta_{ji}
\]
are centered correctly given the earlier histories. Thus the estimator remains
consistent even if the propensity scores are not consistently estimated, provided
the weights are well behaved.

\subsection{Correct $g$-Sequence}

If the treatment mechanisms are consistently estimated, the weighted residual
terms correct bias in the outcome regressions. The intuition is the same as in
the point-treatment AIPTW estimator. The inverse probability weights make the
observed regime-compatible histories representative of the target population,
and the residual terms repair the g-computation plug-in estimate.

For $K=1$, if $g_{0,n}$ and $g_{1,n}$ converge to the true treatment mechanisms,
then the weighted residuals at the two stages remove the drift caused by
incorrect $\Qbar_{1,n}$ and $\Qbar_{2,n}$.

\subsection{Mixed Robustness Patterns Over Time}

With many treatment times, robustness can be more nuanced. For example, an
estimator may be consistent if:
\begin{itemize}
  \item the treatment mechanisms are correct through early time points, and
  \item the outcome regressions are correct from a later time point onward.
\end{itemize}
This occurs because backward recursion only needs enough correct information to
bridge each step between the final outcome and baseline. However, this should
not be interpreted as permission to fit careless models. A misspecified nuisance
function can still harm finite-sample performance, inflate variance, and damage
inference. The valid mixed-robustness patterns are estimator-specific, not
arbitrary choices of some correct $g$ terms and some correct $\Qbar$ terms.

\begin{estimationbox}
The longitudinal AIPTW workflow is:
\begin{enumerate}
  \item Fit all time-specific propensity score models $g_{0,n},\ldots,g_{K,n}$.
  \item Fit the sequential outcome regressions
  $\Qbar_{K+1,n},\ldots,\Qbar_{1,n}$.
  \item Compute the g-computation plug-in estimate
  $n^{-1}\sum_i\Qbar_{1,n}(L_{0i})$.
  \item Add weighted residual corrections at each time point.
  \item Estimate uncertainty using the empirical variance of the estimated
  influence-function contributions, with caution when nuisance fits are highly
  adaptive.
\end{enumerate}
\end{estimationbox}

\begin{keyidea}
Longitudinal AIPTW corrects sequential g-computation by adding weighted residuals
at every treatment time. Its robustness comes from using both the outcome
process and the treatment process.
\end{keyidea}
